\chapter{Giới thiệu}

\section{Bối cảnh và lý do chọn đề tài}

Trong bối cảnh cuộc Cách mạng Công nghiệp 4.0, chuyển đổi số đang trở thành xu thế tất yếu và là yêu cầu cấp thiết đối với các cơ sở giáo dục đại học. Một trường đại học hiện đại không chỉ là nơi đào tạo và nghiên cứu mà còn là một tổ chức vận hành phức tạp với hàng chục ngàn sinh viên, giảng viên và nhân viên. Để quản lý bộ máy này, nhà trường thường sử dụng nhiều hệ thống thông tin chuyên biệt như: Hệ thống quản lý đào tạo, Hệ thống quản lý nhân sự, Hệ thống quản lý tài chính - kế toán, Hệ thống quản lý thư viện và các hệ thống E-learning.

Tuy nhiên, thực trạng phổ biến tại nhiều trường đại học hiện nay là sự phân mảnh dữ liệu. Các phòng ban chức năng thường vận hành các hệ thống phần mềm riêng lẻ, được phát triển bởi các nhà cung cấp khác nhau hoặc vào các thời điểm khác nhau. Điều này dẫn đến những vấn đề nghiêm trọng:

\begin{itemize}
    \item \textbf{Thiếu tính nhất quán dữ liệu:} Thông tin của cùng một sinh viên (ví dụ: trạng thái đóng học phí) có thể sai lệch giữa Phòng Tài chính và Phòng Đào tạo, gây khó khăn cho việc đăng ký môn học hoặc xét tốt nghiệp.
    \item \textbf{Khó khăn trong tích hợp và chia sẻ:} Việc kết nối giữa các hệ thống thường được thực hiện theo mô hình điểm - điểm. Khi số lượng hệ thống tăng lên, mạng lưới kết nối trở nên chằng chịt, khó bảo trì và mở rộng.
    \item \textbf{Khó khăn trong khai thác thông tin:} Dữ liệu không được cung cấp tập trung, khiến việc truy xuất thông tin toàn cảnh về hoạt động của trường gặp nhiều trở ngại.
    \item \textbf{Khó khăn trong báo cáo và ra quyết định:} Lãnh đạo nhà trường thiếu một cái nhìn toàn cảnh về hoạt động của trường do dữ liệu nằm rải rác, không được tổng hợp kịp thời.
\end{itemize}

Trước nhu cầu cấp thiết về việc đồng bộ hóa dữ liệu nội bộ và kết nối với các hệ thống bên ngoài (như CSDL Giáo dục Quốc gia - HEMIS), đề tài \textbf{"Tìm hiểu và hiện thực hệ thống trục tích hợp dữ liệu cho trường đại học"} được lựa chọn thực hiện. Đề tài tập trung giải quyết bài toán phân mảnh dữ liệu bằng cách xây dựng một trục tích hợp dữ liệu hiện đại, đảm bảo tính chính xác, bảo mật và sẵn sàng của dòng chảy dữ liệu trong môi trường giáo dục.

\section{Mục tiêu nghiên cứu}

Mục tiêu tổng quát của đề tài là xây dựng một hệ thống trục tích hợp dữ liệu có khả năng thu thập, chuẩn hóa, lưu trữ và chia sẻ dữ liệu giữa các hệ thống thông tin trong trường đại học một cách tự động và nhất quán.

Các mục tiêu cụ thể bao gồm:

\begin{enumerate}
    \item \textbf{Nghiên cứu giải pháp:} Tìm hiểu các phương pháp tích hợp dữ liệu phổ biến. Nghiên cứu các kỹ thuật về tích hợp dữ liệu, quản trị dữ liệu và các giải pháp đảm bảo chất lượng dữ liệu và bảo mật hệ thống.
    \item \textbf{Thiết kế kiến trúc hệ thống:} Đề xuất một kiến trúc trục tích hợp dữ liệu tối ưu cho ngữ cảnh trường đại học, khắc phục được nhược điểm của các mô hình cũ, đảm bảo tính linh hoạt, khả năng mở rộng và dễ dàng bảo trì.
    \item \textbf{Hiện thực hóa giải pháp:} Xây dựng trục tích hợp dữ liệu với các thành phần cốt lõi:
    \begin{itemize}
        \item \textit{Data Integration Server:} Thu thập và chuẩn hóa dữ liệu từ các nguồn khác nhau.
        \item \textit{Data Sharing Server:} Cung cấp API chuẩn hóa cho các dịch vụ khác trong nội bộ trường học.
        \item \textit{Centralized Database:} Cơ sở dữ liệu dùng chung của trục tích hợp dùng để lưu trữ dữ liệu nhất quán.
        \item \textit{System Management:} Hệ thống cấu hình và quản trị trục tích hợp.
        \item \textit{Các thành phần hỗ trợ:} Định thời, bảo mật API, và giám sát.
    \end{itemize}
    \item \textbf{Đánh giá và kiểm thử:} Kiểm chứng khả năng hoạt động, hiệu năng và tính đúng đắn của dữ liệu sau khi tích hợp.
\end{enumerate}

\section{Phạm vi nghiên cứu}

Để đảm bảo tính khả thi trong khuôn khổ thời gian của đồ án, phạm vi nghiên cứu được giới hạn như sau:

\begin{itemize}
    \item \textbf{Về dữ liệu:} Tập trung vào các nhóm dữ liệu cốt lõi bao gồm: Dữ liệu Viên chức - Người lao động và dữ liệu người học.
    \item \textbf{Về hệ thống nguồn:} Nhóm thực hiện sẽ tự chuẩn bị dữ liệu giả lập để mô phỏng các nguồn dữ liệu đầu vào, do các hệ thống thực tế nằm ngoài phạm vi kiểm soát của đề tài.
    \item \textbf{Về công nghệ:} Sử dụng các công nghệ mã nguồn mở hiện đại phù hợp với kiến trúc Microservices bao gồm NestJS, Next.js, Docker, PostgreSQL, MinIO, ELK Stack.
\end{itemize}

\section{Phương pháp nghiên cứu}

Đề tài áp dụng kết hợp các phương pháp nghiên cứu sau:

\begin{itemize}
    \item \textbf{Phương pháp nghiên cứu tài liệu:} Tổng hợp lý thuyết từ các sách chuyên khảo, bài báo khoa học và tài liệu kỹ thuật về kiến trúc phần mềm, Data Integration, Enterprise Service Bus và API Gateway.
    \item \textbf{Phương pháp phân tích và thiết kế hệ thống:} Áp dụng tư duy thiết kế Domain-Driven để phân tích nghiệp vụ trường đại học. Sử dụng ngôn ngữ UML để mô hình hóa kiến trúc và luồng dữ liệu.
    \item \textbf{Phương pháp thực nghiệm:}
    \begin{itemize}
        \item Xây dựng môi trường giả lập để triển khai hệ thống.
        \item Áp dụng quy trình phát triển phần mềm Agile/Scrum.
        \item Sử dụng CI/CD để tự động hóa quy trình kiểm thử và triển khai.
    \end{itemize}
    \item \textbf{Phương pháp so sánh, đánh giá:} So sánh kiến trúc đề xuất với các mô hình truyền thống để làm rõ ưu/nhược điểm.
\end{itemize}

\section{Đóng góp của luận văn}

Luận văn mang lại những đóng góp chính về cả mặt lý thuyết và thực tiễn:

\begin{itemize}
    \item \textbf{Về mặt kiến trúc:} Đề xuất một mô hình kiến trúc tham chiếu cho trục tích hợp dữ liệu trường đại học dựa trên tư duy \textit{API-led Connectivity}. Kiến trúc này phân tách rõ ràng giữa lớp Hệ thống, lớp Quy trình và lớp Trải nghiệm, giúp tăng tính linh hoạt và khả năng tái sử dụng dịch vụ.
    \item \textbf{Về mặt kỹ thuật:} Cung cấp một bộ giải pháp kỹ thuật cụ thể và quy trình vận hành để hiện thực hóa trục tích hợp, bao gồm các cơ chế nâng cao như lập lịch động, sao lưu và phục hồi điểm thời gian, và bảo mật đa lớp.
    \item \textbf{Về mặt ứng dụng:} Sản phẩm của luận văn có thể được sử dụng làm nền tảng cơ sở để triển khai thực tế, đóng vai trò là nơi cung cấp dữ liệu dùng chung cho các ứng dụng nội bộ trong trường đại học, giải quyết bài toán đồng bộ dữ liệu.
\end{itemize}